\documentstyle[11pt]{article}
\oddsidemargin 0.mm
\evensidemargin 0.mm
\pagestyle{empty}
\parskip=2pt
\setlength{\textwidth}{160.mm}
\setlength{\textheight}{240.mm}
\setlength{\topmargin}{-15mm}
%\headheight 0mm \headsep 0mm
%\footheight 0mm
%\setlength{\footskip}{40pt}
%\def\thepage{}
%
\begin{document}
\begin{center}
{\large {\bf
A Use of Computer Algebra in the Inverse Problem of the
Birkhoff-Gustavson Normalization
}}
\\
\bigskip
%[.5cm]
Yoshio UWANO
\\
\medskip
Department of Applied Mathematics and Physics,
Kyoto University \\
Kyoto 606-8501, Japan
\end{center}
%
\smallskip
\begin{center}
{\bf abstract}
\end{center}
%\medskip
The Birkhoff-Gustavson (BG) normalization has been utilized
very effectively in various studies in Hamiltonian systems.
For example, for a given Hamiltonian system with a semi-simple
stable equilibrium point, the BG normal-form power series
of the given Hamiltonian around the equilibrium point
provides an \lq approximate Hamiltonian system' if truncated
up to a finite degree: the approximate system provides a
good account of the surfaces of section with sufficiently
small energies. The example above implies that finding the class
of Hamiltonian systems sharing a certain polynomial in BG normal
form amounts to finding a class of Hamiltonian systems exhibiting
similar surfaces of sections.
{\it The inverse problem of the BG normalization} has been
thereby posed by the speaker [U1,U2]:
{\it What kind of polynomial Hamiltonians can be brought into
a given polynomial Hamiltonian in the Birkhoff-Gustavson normal form ?}
Since elementary algebraic operations, differentiation, and integration
of polynomials have to be repeated many times both in the inverse and
the direct problems, computer algebra is worth applying to these problems;
the program named \lq $\mbox{ANFER}$' working on REDUCE has been
proposed [U1, U2].
\par
The inverse problem of the BG normalization is expected to
provide rich subjects related with the integrability of Hamiltonian systems.
For example, in solving the inverse problem associated with
the two-dimensional perturbed harmonic oscillators with homogeneous
cubic-polynomial potentials (PHOCP), the Bertrand-Darboux integrability
condition comes out as a necessary and sufficient condition for the PHOCPs
to share a BG normal form up to degree-4 with a certain the two-dimensional
perturbed harmonic oscillators with homogeneous quartic-polynomial
potentials
(PHOQP).
\par
The aim of talk is to present
\begin{description}
\item[{(1)}]
the algorithm of ANFER solving the inverse problem,
\item[{(2)}]
how ANFER works effectively in the inverse problem
associated with the PHOCPs,
\item[{(3)}]
further results in the inverse problem with the PHOCPs
obtained from the viewpoint of separability,
\end{description}
not only from the viewpoint of computer algebra but also
from that of mathematical analysis: computer algebra and mathematical
analysis have been combined very well in the inverse problem, indeed.
\medskip
\baselineskip=14pt
\begin{center}
{\bf References}
\end{center}
\begin{description}
\item[{[U1]}]
Y.Uwano, Journal of Physics {\bf A}, {\bf 33}, 6635--53 (2000).
\item[{[U2]}]
Y.Uwano {\it et.~al.},
{\it Computer Algebra in Scientific Computing},
V.Ganzha {\it et.~al.} eds.,
441--61 (Springer-Verlag, 1999).
\end{description}
\end{document}
